\documentclass{article}
\usepackage[margin=1in]{geometry}
\usepackage{graphicx}
\usepackage{xcolor}
\usepackage{float}
\usepackage{amsmath}
\usepackage{cite}
\usepackage{hyperref}
\usepackage{indentfirst}
\graphicspath{{..} {./images}}

\definecolor{navy-blue}{rgb}{0.22,0.38,0.71}

\renewcommand{\contentsname}{\vspace*{-2\baselineskip}}

\hypersetup{
	colorlinks,
	linkcolor=black,
	urlcolor=blue,
	citecolor=black
}
  		
\begin{document}
\begin{titlepage}
	\centering
	{\huge Lab 8 - SDR: Software Defined Radar}\\[0.25 in]
	\includegraphics[width=0.6\textwidth]{ua_logo.png}\\[0.25 in]
	{\large \textbf{ECE 531 - Software Defined Radio\\[0.25 in]
	April 28, 2025\\[0.25 in]}}
	{\large Owen Sowatzke, osowatzke@arizona.edu\\[0.05 in]
	Department of Electrical \& Computer Engineering\\[0.05 in]
	University of Arizona, Tucson, AZ 85721\\[0.5 in]}
	\hypersetup{linkcolor=navy-blue}
	\noindent\hrulefill
	\tableofcontents
	\noindent\hrulefill
\end{titlepage}

% \setlength{\parindent}{0pt}

\section{Introduction}
%Introduction to the laboratory experiment, including a brief description of the objectives and goals.

\section{Procedure}
% Detailed explanation of the laboratory experiment, including the design, implementation, and testing of the system.

In this lab, we implement a monostatic CW radar with the PlutoSDR. We specifically use GNU radio to implement our radar and test it using a box fan. In our data, we explain the strong DC return and discuss physical changes that can be made to reduce the power of this return. Using the measured doppler shift, we also approximate the angular velocity of the box fan blades. Then, we discuss how an interfering signal could affect the operation of our radar and propose changes to make the design more resilient to interference. Finally, we discuss how the design can be updated to support range measurements.

\section{Results}
% Results and discussion of the laboratory experiment, including captured outputs, observations, and responses to laboratory questions.

Two common causes of the DC return in our data are clutter (from stationary returns) and transmitter leakage. We can minimize clutter with a tighter antenna pattern or by operating in an controlled environment, such as an anechoic chamber. However, it is usually not possible for us to constrain our environment. Antenna isolation achieved through a tighter antenna pattern or other techniques can be help us achieve better antenna isolation. We can also consider other design such as pulsed doppler radars. However, these radars may not be applicable in all cases (especially with close targets). In the lab, we used a notch filter to remove DC. However, we could have also removed or ignored the offending doppler bins in post-processing.

An interfering signal would negatively impact the CW radar operation. If it was a wideband signal (such as 802.11 WiFi), it could mask the entire bandwidth of our signal and make detection more difficult. Conversely, if the interfering signal was narrow (worse case a tone) it could lead to false detections.
 
\section{Conclusion}
% Conclusions to the overall lab that discuss meaningful lessons learned and other takeaways from the assignment. (Important)

\bibliographystyle{IEEEtran}
\bibliography{sources}{}

\end{document}